% part: intuitionistic-logic
% chapter: propositions-as-types
% section: terms-to-proofs

\documentclass[../../../include/open-logic-section]{subfiles}

\begin{document}

\olfileid{int}{pty}{tp}

\olsection{Ngrakit Maneh \usetoken{P}{derivation} saka Terma Bukti}

Saiki kita nimbang arah kosokbaline: ngowahi terma bukti~$N$ dadi
!!a{proof} maneh ing \Log{N2i}. Umume, iki mung bisa ditindakake
relatif marang konteks sing ngandhut pasangan $x : !A$ kanggo saben
variabel~$x$ sing kedadeyan bebas ing terma bukti kasebut. Nanging,
kita wiwiti saka conto tanpa variabel bebas, yaiku terma
\[
\lambd[\typeof{x}{!A \land
  !B}][\pair{\proj{2}{x}}{\proj{1}{x}}]
\]
Terma iki awujud $\lambd[\typeof{x}{!A \land !B}][N_1]$. Amarga
siji-sijine aturan \Log{tN2i} sing ngasilake terma kaya mangkono
iku~\Intro{\lif}, kita ngerti yen !!{proof} sing dikode dening~$N$
kudu dipungkasi kanthi~\Intro{\lif}, yaiku dipungkasi kaya mangkene
  \[
  \Axiom$x : !A\land !B \fCenter \pair{\proj{2}{x}}{\proj{1}{x}} : 
    !C$
  \RightLabel{\Intro\lif}
  \UnaryInf$\fCenter \lambd[\typeof{x}{!A \land
  !B}][\pair{\proj{2}{x}}{\proj{1}{x}}] : (!A \land !B) \lif !C$
  \DisplayProof
  \]
Kita durung ngerti $!C$, nanging kita ngerti yen antesedhene iku
$!A \land !B$, lan konteks premise kudu ngandhut $x:!A \land !B$.

Saiki terma bukti sing digandhengake karo premis wis ditemtokake:
$\pair{\proj{2}{x}}{\proj{1}{x}}$. !!{proof} sing \emph{dikode dening
terma iki} kudu dipungkasi kanthi \Intro{\land}. Kita isih durung
ngerti $!C$, nanging amarga iku konklusi saka~\Intro{\land}, kita
gerti yen kudu awujud konjungsi, yaiku $!C \ident (!C_1 \land !C_2)$.
Kita nerusake pangrakitan maneh:
  \[
  \Axiom$x : !A\land !B \fCenter \proj{2}{x} : !C_1$
  \Axiom$x : !A\land !B \fCenter \proj{1}{x} : !C_2$
  \RightLabel{\Intro\land}
  \BinaryInf$x : !A\land !B \fCenter \pair{\proj{2}{x}}{\proj{1}{x}} : 
    !C_1 \land !C_2$
  \RightLabel{\Intro\lif}
  \UnaryInf$\fCenter \lambd[\typeof{x}{!A \land
  !B}][\pair{\proj{2}{x}}{\proj{1}{x}}] : (!A \land !B) \lif (!C_1 \land !C_2)$
  \DisplayProof
  \]
Terma bukti $\pi_2(x)$ nuduhake yen sekuen paling ndhuwur ing kiwa
dipikolehi kanthi \Elim{\land}[2], yaiku premise awujud
$!D \land !C_1$. Ing tengen, kita bisa nemtokake yen inferensi sing
ngasilake $!C_2$ kudu \Elim{\land}[1], lan premise awujud
$!C_2 \land !E$.
  \[
  \Axiom$x : !A\land !B \fCenter x: !D\land !C_1$
  \RightLabel{\Elim\land[2]}
  \UnaryInf$x : !A\land !B \fCenter \proj{2}{x} : !C_1$
  \Axiom$x : !A\land !B \fCenter x: !C_2\land !E$
  \RightLabel{\Elim\land[1]}
  \UnaryInf$x : !A\land !B \fCenter \proj{1}{x} : !C_2$
  \RightLabel{\Intro\land}
  \BinaryInf$x : !A\land !B \fCenter \pair{\proj{2}{x}}{\proj{1}{x}} : 
    !C_1 \land !C_2$
  \RightLabel{\Intro\lif}
  \UnaryInf$\fCenter \lambd[\typeof{x}{!A \land
  !B}][\pair{\proj{2}{x}}{\proj{1}{x}}] : (!A \land !B) \lif (!C_1 \land !C_2)$
  \DisplayProof
  \]
Amarga terma bukti saka sekuen paling ndhuwur saiki awujud variabel,
kita ngerti yen sekuen kasebut kudu aksioma. Iki nemtokake $!C_1$,
$!C_2$, $!D$, lan $!E$: $!D \ident !A$ lan $!C_1 \ident !B$, amarga
yen ora, sekuen kiwa dudu aksioma; lan $!C_2 \ident !A$ lan
$!E \ident !B$, amarga yen ora, sekuen tengen dudu aksioma. Kita wis
ngrakit maneh bukti kanthi jangkep:
  \[
  \Axiom$x : !A\land !B \fCenter x: !A\land !B$
  \RightLabel{\Elim\land[2]}
  \UnaryInf$x : !A\land !B \fCenter \proj{2}{x} : !B$
  \Axiom$x : !A\land !B \fCenter x: !A\land !B$
  \RightLabel{\Elim\land[1]}
  \UnaryInf$x : !A\land !B \fCenter \proj{1}{x} : !A$
  \RightLabel{\Intro\land}
  \BinaryInf$x : !A\land !B \fCenter \pair{\proj{2}{x}}{\proj{1}{x}} : 
    !B \land !A$
  \RightLabel{\Intro\lif}
  \UnaryInf$\fCenter \lambd[\typeof{x}{!A \land
  !B}][\pair{\proj{2}{x}}{\proj{1}{x}}] : (!A \land !B) \lif (!B \land !A)$
  \DisplayProof
  \]


\end{document}
